\chapter{Beitrag von Sascha Lauk}\label{ch:sascha}

\saschaPH{Dieses Kapitel schreibt Sascha Lauk selbst. Die folgende Gliederung ist nur ein Vorschlag, damit Inhalt und Umfang zum restlichen Bericht passen; sie kann frei geändert werden.}

\section{Aufgabenbereich und Einordnung}
\saschaPH{Welche Teile des Systems hat Sascha bearbeitet (z.\,B.\ Kapitelbezug zu Architektur, Crawler, Klassifikation, Evaluation, Frontend)? Wie ist die Schnittstelle zu den Arbeiten in den Kapiteln~\ref{ch:architektur} bis~\ref{ch:klassifikation}?}

\section{Vorgehen und Entscheidungen}
\saschaPH{Eingeschlagene Wege, Alternativen, Sackgassen, wie in Kapitel~\ref{ch:diskussion} für den Hauptteil.}

\section{Ergebnisse und Kennzahlen}
\saschaPH{Messbare Ergebnisse des eigenen Teils (Zahlen, Diagramme, Tabellen).}

\section{Probleme und Lehren}
\saschaPH{Probleme im Betrieb oder in der Entwicklung, Lehren.}

\section{Wochenübersicht}
Die wochenweise Aufschlüsselung von Sascha Lauk steht in Tabelle~\ref{tab:wochen-sascha}.
